\section{Step 34: F24 Role-Obstruction Discriminator}

\paragraph{Deflationary result.}
F24 role-obstruction does not distinguish the reference support.  The non-tuned predicate \(O_s=0\) is a principled descent test, but the reference has \(O_s=1\), so it fails the F24-descending cut.

\paragraph{Instantiation.}
The access quotient \(q\) is the canonical non-abelian representation role.  The role readout \(s\) is the abelian charge role.  The obstruction \(O_s\) counts same-\(q\) pairs with different \(s\).  The discriminator \(O_s=0\) asks whether the charge role descends through the non-abelian access quotient.

\paragraph{Computation.}
The corrected Step-33 final carrier is reproduced with \(80\) survivors.  The obstruction distribution is \(O_s=0:10\), \(O_s=1:34\), and \(O_s=2:36\).  The reference has \(O_s=1\), a value shared by \(34\) survivors.

\paragraph{Verdict.}
The result is \texttt{TYPED\_NO\_GO\_TYPE\_LIMIT}.  F24 supplies a real can-fail discriminator, but not one that selects the reference.  The next grammar delta is a different kind of principle: higher-layer descent/content-cascade criterion or observed-input boundary.

\paragraph{Boundary.}
No physical gauge-structure theorem, physical value, real selection mechanism, or frame-transfer status upgrade is claimed.
